\subsection{Cyclic support on three and five objects at eighteen arrows}

\begin{lemma}[Proved: a binary collage into the five-element band monoid]
\label{CP2:A-P5-binary-source-block}
Let a finite category be a collage of a source category $S$ and
a sink corner $P_5$, the identity-adjoined $2\times2$ rectangular
band. Suppose every stalk $F(x)=C(x,*)$, for $x\in S$, is nonempty
and has at most two elements. If $H^2(BS;\mathbb Q)=0$, all rational
degree-two products on $BC$ vanish.
\end{lemma}
\begin{proof}
Write the band elements as $e_{ij}$, with
$e_{ij}e_{kl}=e_{il}$. If any stalk contains a point fixed by all
of $P_5$, that arrow gives a natural transformation from the
constant functor at its source object to the sink-corner inclusion.
The inclusion $BP_5\to BC$ is nullhomotopic. Every global second
class therefore restricts to zero on both blocks, and
Lemma~\ref{CP2:A-two-block} kills every product.

Assume no stalk has a $P_5$-fixed point. Every stalk then has two
elements. An idempotent transformation of a two-element set is
either the identity or a constant. If one band element acts as
the identity, the rectangular-band identities force all of them
to do so: in its row and column their products with that identity
give the identity, so those idempotents are bijective; then all
remaining entries are their products. This would give fixed
points. Thus every band element acts constantly. The relation
$e_{ij}e_{kl}=e_{il}$ shows that the constant value depends only
on $i$. The two row values must be distinct, since otherwise their
common value is fixed. Hence each stalk is canonically the
two-element row act $U=P_5e$ for a minimal idempotent $e$.

Every map between two stalks commutes with these two row constants,
and must preserve both row values. Under the specified
identifications all source-category arrows therefore act by the
identity on $U$. The cross-arrow bar is consequently
$B(P_5,U)\times BS$, and its projection to $BS$ is a homotopy
equivalence. Indeed $U=P_5e$ is an equivariant retract of the
regular left act $P_5$, by inclusion and $m\mapsto me$.
The action category of the regular left act has initial object
$1$, so its classifying space is contractible; its retract
$B(P_5,U)$ is contractible as well.

The actual collage homotopy pushout now identifies $BC\simeq BP_5$,
because the cross-arrow bar maps equivalently to the source block.
The latter space has type $S^2$, by
Corollary~\ref{CP2:A-five-types}, and its degree-two products are
zero. This proves the remaining case.
\end{proof}

\begin{theorem}[Proved: cyclic three-object support through eighteen]
\label{CP2:A-strong-three-eighteen}
A finite category on three objects with cyclic support and at most
eighteen arrows cannot have classifying space of type $\CP^2$.
\end{theorem}
\begin{proof}
Assume a connected counterexample. An object-retract deletion
saves at least three arrows and leaves at most fifteen; the
complete sixteen-arrow baseline excludes it. Thus there is no
object retract. If all three objects form a strong component,
Lemma~\ref{CP2:A-high-Rees} gives finite-group type unless both
total minimal row and column counts are at least four. In that
case the three identities are outside the ideal and there are
at least $4\cdot4+3=19$ arrows, a contradiction.

There remains a proper strong pair $S$ and one corner monoid $N$.
By opposition take $S$ to be the source component. Write
$q=|S|$, $m=|N|$, and $h$ for the total outgoing stalk cardinality.
Both stalks are nonempty, so
\[
 q+m+h\le18,\qquad q\ge6,\qquad h\ge2.
\]
If $h=2$, the full collage is a categorical join and has zero
positive rational products. Hence assume $h\ge3$. If both
$H^2(BS;\mathbb Q)$ and $H^2(BN;\mathbb Q)$ are zero, the
two-block lemma applies.

Suppose first that $H^2(BS;\mathbb Q)\ne0$. Its total minimal
row and column counts $r,c$ are both at least three; otherwise
one local count is one at each object and $S$ has finite-group
type. Since $q\le14$, the minimal group is trivial and
$rc+2\le q$ leaves only
\[
 (r,c)=(3,3),(3,4),(4,3).
\]
If $m\le2$ and each local row count is at most two, the trivial
right rational $N$-module is projective and
Corollary~\ref{CP2:A-baseline-two-row-instances} annihilates the
collage's second products. The only remaining local-row
distribution is $(1,3)$, in either order, with $(r,c)=(4,3)$.
It forces $q=14$, so the internal category consists exactly of
its twelve-element minimal band and its two identities.
At its one-row object $b$, a local minimal idempotent fixes every
incoming arrow from the other object, and the only unit is the
identity. There are no external entrances into $S$. Minimal cross
arrows supply the factorization of this idempotent required in
Lemma~\ref{CP2:A-high-one-sided}. Its one-sided object deletion
is an actual homotopy equivalence. It removes at least the two
corner arrows and two internal cross arrows, leaving at most
fourteen arrows, which the complete baseline excludes.

If $m\ge3$, then $q\le12$, so necessarily $r=c=3$. For $q=11$
the pair is again exactly its minimal band and its identities,
and the same one-row object deletion applies. For $q=12$ the
budget forces $m=3,h=3$. Its two outgoing stalks have sizes one
and two. Lemma~\ref{CP2:A-high-small-binary}, with source block
$S$ and the order-three sink, makes the square of its single
global second class zero. This exhausts the case of nonzero
second cohomology on $S$.

Finally suppose $H^2(BS;\mathbb Q)=0$ and
$H^2(BN;\mathbb Q)\ne0$. Then $m\ge5$. A singleton outgoing
stalk is fixed by all $N$ actors and the fixed-arrow corollary
applies. Therefore both stalks have size at least two, giving
$h\ge4$ and $q\le9$. The minimal group of $S$ is trivial:
a nontrivial group would already require
$2\cdot2\cdot2+2=10$ internal arrows. Thus
$rc+2\le9$ leaves either $r=2$, or $r=3,c=2$.

When $r=2$, the source-component collapse in
Lemma~\ref{CP2:A-high-one-sided} replaces $S$ by a trivial
object with the actual outgoing orbit act. The resulting space
is an ordinary act cone over $N$, with $m\le8$ from
$q\ge6,h\ge4$. Its homology is bounded, and the bounded-eleven
act theorem excludes its nonzero square. If $r=3,c=2,q=8$,
the pair is its six-element band plus identities; the one-row
object deletion above again applies.

The only new numerical case is $q=9,m=5,h=4$. The order-five
corner with nonzero second cohomology is necessarily $P_5$,
and the two stalks both have size two. Lemma~\ref{CP2:A-P5-binary-source-block}
applies since $H^2(BS;\mathbb Q)=0$. Its zero-product conclusion
contradicts the projective-plane square. All cyclic-support
possibilities have been excluded.
\end{proof}

\begin{corollary}[Proved: cyclic five-object support through eighteen]
\label{CP2:A-strong-five-eighteen}
A finite category on five objects with cyclic support and at
most eighteen arrows cannot have classifying space of type $\CP^2$.
\end{corollary}
\begin{proof}
Here $n\le18=2(5)+8$. Corollary~\ref{CP2:A-high-nine-support}
states that a projective-plane candidate in this range has acyclic
support. This contradicts the hypothesis directly.
\end{proof}
